
\section*{\underline{Exposé XIV}\\[3pt] \underline{\underline{PROBLÈMES OUVERTS EN THÉORIE DES INTERSECTIONS}}}
\begin{center}
\underline{par A. Grothendieck}
\end{center}

Dans le présent exposé nous passons en revue quelques uns des problèmes liés aux questions rencontrées dans ce séminaire, et qui mériteraient d'être étudiés. Nous laissons de côté les problèmes liés à la \uline{théorie locale} des intersections à la Serre [14], que nous n'avons pas eu à utiliser dans le Séminaire, et renvoyons notamment à loc. cit. pour les deux conjectures fondamentales qui restent ouvertes dans cette direction (loc. cit. Chap V \S~4). Nous passons sous silence aussi les problèmes, fort importants pour les applications arithmétiques, liés à la \uline{structure des groupes de classes de cycles}, sur un schéma algébrique projectif et lisse notamment; nous renvoyons pour la discussion de ces problèmes à S. KLEIMAN, Algebraic cycles and the Weil conjectures dans Dix exposés sur la cohomologie des schémas (Masson-North Holland Publishing Cie). Dans le n° 4 du présent exposé, nous donnons également quelques compléments au Séminaire, concernant notamment l'\uline{anneau de Chow}, qui n'avaient trouvé leur place dans les exposés précédents.

\section*{1. Opérations $\Lambda^i$ dans la catégorie dérivée $D(X)$\SGAAuthorFootMark{*}}
\SGAAuthorFootText{*}{cf. note de la page 23.}

Soit $X$ un topos localement annelé, et soit $K^\bullet(X)$ l'anneau des classes de complexes parfaits sur $X$ (Exp. IV). On aimerait définir dans $K^\bullet(X)$ une structure de $\lambda$-anneau (Exp. V), fonctorielle en $X$, telle que l'homomorphisme naturel
\[
 K^\bullet_{\mathrm{naif}}(X)\longrightarrow K^\bullet(X)
\]
(où le premier membre est l'anneau des classes de Modules localement libres sur $X$, muni de la $\lambda$-structure définie dans Exp. V) soit un -homomorphisme Dans ce but, il y aurait lieu de définir des foncteurs (non addifs en général de ``puissance extérieure)
\[
 \Lambda^i:D^-(X)\longrightarrow D^-(X),
\]
transformant complexes pseudo-cohérents (resp. parfaits) en complexes pseudo-cohérents (resp. parfaits), et ayant les propriétés formelles nécessaires, vis à vis des triangles exacts, pour définir une $\lambda$-structure sur $K^\bullet(X)$ par la formule
\begin{equation}\tag{1.1}
 \lambda^i(\mathrm{c}\ell^\bullet(L_\bullet))=\mathrm{c}\ell^\bullet(\Lambda^i L_\bullet).
\end{equation}

